\section*{अध्याय \ref{ch:b2:liquid-vapour-diagrams} --- द्वि-घटकी द्रव--वाष्प आरेख और आसवन}

\begin{solution}{exo:b2:liquid-vapour-diagrams:1}
बेन्ज़ीन \qty{80.1}{\degreeCelsius}, टॉलूईन \qty{110.6}{\degreeCelsius}। संघटन 0.2
वाला द्रव लगभग \qty{102}{\degreeCelsius} पर उबलना आरंभ करता है; उसकी पहली
वाष्प का संघटन 0.38 है।
\end{solution}

\begin{solution}{exo:b2:liquid-vapour-diagrams:2}
$p = 0.5 \times 1.010 + 0.5 \times 0.388 = \qty{0.699}{bar}$;
$y = 0.505/0.699 = 0.722$।
\end{solution}

\begin{solution}{exo:b2:liquid-vapour-diagrams:3}
$n_V/n = (0.30 - 0.20)/(0.45 - 0.20) = 0.40$: वाष्प के \qty{4.0}{mol} और
द्रव के \qty{6.0}{mol}। जाँच: $6.0 \times 0.20 + 4.0 \times 0.45 = 3.0 = 10
\times 0.30$।
\end{solution}

\begin{solution}{exo:b2:liquid-vapour-diagrams:4}
(a) $v = 2 + 2 - 1 = 3$: $T$, $p$ और $x$ स्वतंत्र हैं। (b) $v = 2 + 2 - 2 = 2$;
नियत दाब पर 1। (c) \omterm{def:b2:liquid-vapour-diagrams:azeotrope}{स्थिरक्वाथी} संबंध $x = y$ जोड़ता है: $v = 1$; नियत
दाब पर 0 --- \omterm{def:b2:liquid-vapour-diagrams:azeotrope}{स्थिरक्वाथी} शुद्ध द्रव की तरह नियत ताप पर उबलता है,
यद्यपि उसका संघटन दाब के साथ बदलता है।
\end{solution}

\begin{solution}{exo:b2:liquid-vapour-diagrams:5}
\omterm{def:b2:liquid-vapour-diagrams:bubble-dew}{ओसांक} पर $yp/p_1^* + (1 - y)p/p_2^* = 1$। \qty{100}{\degreeCelsius} पर:
$0.3 \times 1.013/1.80 + 0.7 \times 1.013/0.742 = 1.124 > 1$ (बहुत ठंडा);
\qty{105}{\degreeCelsius} पर: $0.148 + 0.824 = 0.971 < 1$। अंतर्वेशन से $T \approx
100 + 5 \times 0.124/0.153 = \qty{104}{\degreeCelsius}$। पहली बूँद का $x =
yp/p_1^* \approx 0.3 \times 1.013/2.0 = 0.15$ है।
\end{solution}

\begin{solution}{exo:b2:liquid-vapour-diagrams:6}
पहली बूँदों का संघटन उबलते द्रव के ऊपर की वाष्प का होता है,
$y = 0.71$, \qty{92.1}{\degreeCelsius} पर। चूँकि बेन्ज़ीन वरीयता से निकलती है,
द्रव उसमें निर्धन होता जाता है, उसका क्वथनांक बढ़ता है और आसुत भी
निर्धन होता जाता है; अंतिम बूँदें लगभग शुद्ध टॉलूईन होती हैं। सरल आसवन एक
अकेला साम्य चरण है, जो लगातार बदलते द्रव पर दोहराया जाता है: हर
अंश एक मिश्रण है।
\end{solution}

\begin{solution}{exo:b2:liquid-vapour-diagrams:7}
फ़ीड 0.05, \omterm{def:b2:liquid-vapour-diagrams:azeotrope}{स्थिरक्वाथी} के जल वाले पक्ष पर: शीर्ष \omterm{def:b2:liquid-vapour-diagrams:azeotrope}{स्थिरक्वाथी} देता है
(मोल में 0.88, द्रव्यमान में 95~\%), क्वथित्र जल रखता है। फ़ीड 0.95,
एथेनॉल वाले पक्ष पर: शीर्ष अब भी \omterm{def:b2:liquid-vapour-diagrams:azeotrope}{स्थिरक्वाथी}, न्यूनतम क्वथनांक, देता है,
और क्वथित्र शुद्ध एथेनॉल रखता है।
\end{solution}

\begin{solution}{exo:b2:liquid-vapour-diagrams:8}
मिश्रण तब उबलता है जब $0.10 + p^*_{\text{जल}} = 1.013$, अर्थात्
$p^*_{\text{जल}} = \qty{0.91}{bar}$, लगभग \qty{97}{\degreeCelsius} पर। द्रव्यमान अनुपात
$0.10 \times 136/(0.91 \times 18.02) = 0.83$: प्रति ग्राम जल \qty{0.83}{g}
सुगंध।
\end{solution}

\begin{solution}{exo:b2:liquid-vapour-diagrams:9}
कमरे के ताप पर दो परतें L$_2$ और L$_1$ (संघटन 0.4 अंतराल में
है)। तापन पर परतें \omterm{def:b2:liquid-vapour-diagrams:miscibility-gap}{विषम-स्थिरक्वाथी} ताप पर साथ-साथ उबलती हैं,
जो स्थिर रहता है; पहली वाष्प का संघटन \omterm{def:b2:liquid-vapour-diagrams:miscibility-gap}{विषम-स्थिरक्वाथी}
संघटन (बिंदु) है। चूँकि 0.4 बिंदु के L$_2$ वाले पक्ष पर है, परत
L$_1$ पहले उपभुक्त होती है; तब ताप बढ़ता है, शेष L$_2$ अपने \omterm{def:b2:liquid-vapour-diagrams:bubble-dew}{बुलबुला वक्र} के
अनुदिश 1 में निर्धन होती जाती है और वाष्प \omterm{def:b2:liquid-vapour-diagrams:bubble-dew}{ओस वक्र} का अनुसरण करती है, जब तक
बिंदु \omterm{def:b2:liquid-vapour-diagrams:bubble-dew}{ओस वक्र} तक नहीं पहुँचता: अंतिम बूँद वाष्पित हो जाती है।
\end{solution}

\begin{solution}{exo:b2:liquid-vapour-diagrams:10}
$N_{\min} = \ln(99 \times 99)/\ln 2.47 = 9.19/0.904 = 10.2$, जबकि 6.5: दोनों सिरों पर
95~\% से 99~\% शुद्धता तक जाने में आधी से अधिक प्लेटें और
लगती हैं। $x/(1 - x)$ पर हर अतिरिक्त गुणक उतनी ही प्लेटें लेता है, और
शुद्धता शेष अशुद्धि से मापी जाती है, जो केवल गुणोत्तर क्रम में घटती है।
\end{solution}

\begin{solution}{exo:b2:liquid-vapour-diagrams:11}
केवल \omterm{def:b2:liquid-vapour-diagrams:binary-diagram}{समदाबी आरेख} इसे नहीं देता: \qty{80}{\degreeCelsius} पर $p_1^*$ और $p_2^*$
चाहिए, एंटोइन समीकरणों से (1.010 और
\qty{0.388}{bar})। तब $p = p_2^* + x(p_1^* - p_2^*)$ और $p = 1/(y/p_1^* + (1 -
y)/p_2^*)$। द्रव उच्च दाब पर स्थायी है (वाष्प को संपीडित करने से वह
संघनित होती है), \omterm{def:b2:liquid-vapour-diagrams:binary-diagram}{समदाबी आरेख} पर निम्न ताप पर: दोनों आरेख
एक-दूसरे के सापेक्ष उलटे हैं।
\end{solution}

\begin{solution}{exo:b2:liquid-vapour-diagrams:12}
चूँकि $y$, $x$ के साथ बढ़ता है और $y(1 - y) > 0$, $dp/dx$ का चिह्न $y - x$ का है:
दाब $x$ के साथ वहाँ बढ़ता है जहाँ वाष्प 1 में अधिक समृद्ध है। \omterm{def:b2:chemical-potential:ideal-mixture}{आदर्श
मिश्रण} के लिए $y > x$, हर $0 < x < 1$ पर (\cref{prop:b2:liquid-vapour-diagrams:richer}):
$p$ दृढ़तः एकदिष्ट है, उसका कोई चरम नहीं, और कोई \omterm{def:b2:liquid-vapour-diagrams:azeotrope}{स्थिरक्वाथी} नहीं।
\end{solution}

\begin{solution}{pb:b2:liquid-vapour-diagrams:1}
\textbf{1.} $T = B/(A - \log_{10}1.013) - C$: बेन्ज़ीन $1203.835/4.01253 + 53.226 =
\qty{353.2}{K}$, \qty{80.1}{\degreeCelsius}; टॉलूईन $1343.943/4.07266 + 53.773 =
\qty{383.8}{K}$, \qty{110.6}{\degreeCelsius}।
\textbf{2.} \qty{363.15}{K} पर: $p_1^* = \qty{1.361}{bar}$, $p_2^* = \qty{0.542}{bar}$।
\textbf{3.} $x = (1.013 - 0.542)/(1.361 - 0.542) = 0.575$; $y = 0.575 \times
1.361/1.013 = 0.773$।
\textbf{4.} $x = (1.013 - 0.742)/(1.800 - 0.742) = 0.256$; $y = 0.256 \times
1.800/1.013 = 0.455$।
\textbf{5.} \omterm{def:b2:liquid-vapour-diagrams:bubble-dew}{बुलबुला बिंदु} (0, 110.6), (0.256, 100), (0.575, 90), (1, 80.1); ओस
बिंदु (0, 110.6), (0.455, 100), (0.773, 90), (1, 80.1) --- चित्र का
लेंस-आकार।
\textbf{6.} $\frac12(p_1^* + p_2^*) = 1.013$: \qty{90}{\degreeCelsius} पर 0.952,
\qty{95}{\degreeCelsius} पर 1.102; अंतर्वेशन से $90 + 5 \times 0.061/0.150 \approx
\qty{92}{\degreeCelsius}$ (यथातथ \qty{92.1}{\degreeCelsius})।
\textbf{7.} $v = 2$, नियत दाब पर 1: हर ताप $x$ और $y$ नियत करता है,
इसीलिए आरेख दो वक्रों से बना है।
\textbf{8.} $x = (1.013 - 0.636)/(1.569 - 0.636) = 0.404$; $y = 0.404 \times
1.569/1.013 = 0.626$।
\textbf{9.} $n_V = 100 \times (0.500 - 0.404)/(0.626 - 0.404) = \qty{43}{mol}$,
$n_L = \qty{57}{mol}$।
\textbf{10.} वाष्प $43 \times 0.626 = \qty{26.9}{mol}$; द्रव $57 \times 0.404
= \qty{23.0}{mol}$; कुल \qty{49.9}{mol} $\approx 50$।
\textbf{11.} बेन्ज़ीन का $26.9/50 = 54~\%$, ऐसी वाष्प में जिसमें केवल 63~\% है।
\textbf{12.} \qty{100}{\degreeCelsius} पर साम्य वाष्प का $y =
0.455 < 0.5$ है: फ़ीड बिंदु \omterm{def:b2:liquid-vapour-diagrams:bubble-dew}{ओस वक्र} से ऊपर है, फ़ीड पूरा
वाष्प है और कुछ भी पृथक नहीं होता।
\textbf{13.} अपने \omterm{def:b2:liquid-vapour-diagrams:bubble-dew}{ओसांक} पर, \qty{98.8}{\degreeCelsius} (\cref{exo:b2:liquid-vapour-diagrams:5} की
विधि से)।
\textbf{14.} $\alpha = p_1^*/p_2^*$; \qty{80.1}{\degreeCelsius} पर $1.013/0.390 =
2.60$।
\textbf{15.} \qty{110.6}{\degreeCelsius} पर $2.377/1.013 = 2.35$।
\textbf{16.} $\sqrt{2.60 \times 2.35} = 2.47$।
\textbf{17.} $y = xp_1^*/p$ और $1 - y = (1 - x)p_2^*/p$; भाग दीजिए।
\textbf{18.} कुछ भी नहीं निकाला जाता: दो प्लेटों के बीच ऊपर जाती वाष्प और
नीचे आता द्रव द्रव्य की और हर घटक की समान मात्रा ले जाते हैं,
$V y_k = L x_{k+1}$, जहाँ $V = L$, अतः $x_{k+1} = y_k$।
\textbf{19.} $r = x/(1 - x)$ के साथ: $r(y_k) = \alpha\,r(x_k)$ और $x_{k+1} = y_k$, अतः
$r(x_{k+1}) = \alpha\,r(x_k)$ और $r(x_D) = \alpha^N r(x_B)$; $N = \ln[r(x_D)/
r(x_B)]/\ln\alpha$।
\textbf{20.} $N_{\min} = \ln(19 \times 19)/\ln 2.47 = 5.889/0.904 = 6.5$।
\textbf{21.} स्तंभ को उत्पाद देना होता है, इसलिए वह परिमित \omterm{def:b2:liquid-vapour-diagrams:fractional-distillation}{पश्चवाह अनुपात} पर चलता है,
जिसके लिए अधिक प्लेटें चाहिए; और वास्तविक प्लेट साम्य तक नहीं पहुँचती (उसकी
दक्षता एक से कम है)।
\textbf{22.} सात पायदान, अंतिम 0.034 पर समाप्त, 0.05 से नीचे: सतत गिनती में 6.5
प्लेटें, सात पूर्ण प्लेटें (पुनःक्वथित्र को एक गिनकर)।
\textbf{23.} $x = (95/46.07)/(95/46.07 + 5/18.02) = 0.881$।
\textbf{24.} शीर्ष पर अधिक से अधिक \omterm{def:b2:liquid-vapour-diagrams:azeotrope}{स्थिरक्वाथी} (मोल में 0.88); तली पर
जल। पचास प्लेटें \omterm{def:b2:liquid-vapour-diagrams:azeotrope}{स्थिरक्वाथी} को पार नहीं करतीं।
\textbf{25.} पूर्ण पश्चवाह पर $\boldsymbol{N_{\min} \approx 6.5}$ सैद्धांतिक
प्लेटें।
\end{solution}
